% titel: Blatt 7, A2: C2H2, C2H4, H2CCCH2, C4H4 - Lewis, Hybridisierung, Winkel, Ebene
% Teil II: der durchdachte Weg.

\subsection*{Gegeben und gesucht}
Vier Kohlenwasserstoffe $\mathrm{HCCH}$, $\mathrm{H_2CCH_2}$,
$\mathrm{H_2CCCH_2}$ und $\mathrm{H_2CCCCH_2}$. F\"ur jedes: Lewis-Formel,
Hybridisierung der C-Atome, Bindungstypen zwischen C, Bindungswinkel und
die Frage, ob alle H-Atome in einer Ebene liegen.

\subsection*{Ansatz}
Doppelbindung $=$ 1$\sigma$+1$\pi$; Dreifachbindung $=$ 1$\sigma$+2$\pi$;
Einfach $=$ 1$\sigma$. Hybridisierung aus der Zahl der Nachbarn und
$\pi$-Bindungen.

\subsection*{Rechnung}
\begin{enumerate}
  \item[a)] \textbf{Lewis:}
  \begin{itemize}
    \item $\mathrm{HCCH}$ ($\mathrm{C_2H_2}$): $\mathrm{H\equiv C \mathbin{-} C\equiv H}$,
          eine C$\equiv$C-Dreifachbindung.
    \item $\mathrm{H_2CCH_2}$ ($\mathrm{C_2H_4}$): $\mathrm{H_2C=C H_2}$,
          eine C=C-Doppelbindung.
    \item $\mathrm{H_2CCCH_2}$ ($\mathrm{C_3H_4}$, Allen):
          $\mathrm{H_2C=C=C H_2}$ (Zwei Doppelbindungen C=C=C).
    \item $\mathrm{H_2CCCCH_2}$ ($\mathrm{C_4H_4}$, Butatrien):
          $\mathrm{H_2C=C=C=CH_2}$ (drei C=C).
  \end{itemize}
  \item[b)] \textbf{Hybridisierung:}
  \begin{itemize}
    \item $\mathrm{HCCH}$: beide C \textbf{sp} (1 Nachbar C + H,
          Dreifachbindung $\rightarrow$ 2 $\pi$-Bindungen).
    \item $\mathrm{H_2CCH_2}$: beide C \textbf{sp$^2$} (3 Nachbarn,
          1 $\pi$-Bindung).
    \item $\mathrm{H_2CCCH_2}$: End-C \textbf{sp$^2$} (3 Nachbarn),
          Mittel-C \textbf{sp} (2 Nachbarn, 2 $\pi$-Bindungen).
    \item $\mathrm{H_2CCCCH_2}$: End-C \textbf{sp$^2$}, die zwei inneren
          C \textbf{sp}.
  \end{itemize}
  \item[c)] \textbf{Bindungstypen:}
  \begin{center}
  \begin{tabular}{ll}
    Molekül & C--C-Bindung \\ \hline
    $\mathrm{HCCH}$      & 1$\sigma$+2$\pi$ (Dreifach) \\
    $\mathrm{H_2CCH_2}$  & 1$\sigma$+1$\pi$ (Doppel) \\
    $\mathrm{H_2CCCH_2}$ & zwei C=C: je 1$\sigma$+1$\pi$ \\
    $\mathrm{H_2CCCCH_2}$& drei C=C: je 1$\sigma$+1$\pi$ \\
  \end{tabular}
  \end{center}
  \item[d)] \textbf{Bindungswinkel:}
  \begin{itemize}
    \item am sp$^2$-C (Enden): $\approx 120^\circ$ (H--C--H);
    \item am sp-C (linear): $180^\circ$ (C--C--C bzw.\ C--C--H).
    \item $\mathrm{HCCH}$: $180^\circ$ (alle Winkel linear).
  \end{itemize}
  \item[e)] \textbf{Ebene aller H-Atome:}
  \begin{itemize}
    \item $\mathrm{HCCH}$: \textbf{ja} (lineares Molekül, H in einer Linie).
    \item $\mathrm{H_2CCH_2}$: \textbf{ja} (trigonal, $\sigma$-Gerüst
          planar).
    \item $\mathrm{H_2CCCH_2}$ (Allen): \textbf{nein}! Am sp-Zentral-C
          stehen die beiden $\pi$-Systeme senkrecht; die beiden
          $\mathrm{CH_2}$-Ebenen sind gegeneinander um $90^\circ$ verdreht.
    \item $\mathrm{H_2CCCCH_2}$ (Butatrien): \textbf{ja} (die beiden
          End-$\mathrm{CH_2}$-Ebenen liegen in einer gemeinsamen Ebene, die
          $\pi$-Systeme sind gekreuzt aber die Gerüste koplanar).
  \end{itemize}
\end{enumerate}

\begin{ergebnis}
End-C-Atome sind stets sp$^2$ (Winkel $120^\circ$), mittlere C mit zwei
Doppelbindungen sind sp (Winkel $180^\circ$). Die $\pi$-Systeme liegen
senkrecht zur $\sigma$-Bindungsebene. Nur \textbf{Allen}
($\mathrm{H_2CCCH_2}$) hat nicht alle H in einer Ebene.
\end{ergebnis}

\begin{kontrolle}
Konsistenz: Mehr Doppelbindungen $\rightarrow$ mehr sp-C und gr\"ossere
Linearität. Nicht-Planarität des Allens ist die bekannte Sonderregel:
genau die kumulierten Doppelbindungen (C=C=C) verdrehen die
Endgruppen-Ebenen gegeneinander. Butatrien mit gerader Zahl an
Doppelbindungen ist wieder planar.
\end{kontrolle}

\begin{merksatz}
sp$^3$ $= 109{,}5^\circ$ (Einfach), sp$^2$ $= 120^\circ$ (Doppel), sp $=
180^\circ$ (Dreifach/Doppel). Bei kumulierten Doppelbindungen (C=C=C)
steht die Zahl der $\pi$-Systeme; eine ungerade Zahl verdreht die
End-Flügel um $90^\circ$ aus der Ebene.
\end{merksatz}
